\scanpage{56}
% archive foliation: 55
$\Sigma^*_{\underline J}$ défini \deleted{\unsure{par}} à partir de la \deleted{fibré} conique \deleted{\unsure{def}} d'équation homogène
\[
(12)\qquad\sum_{i\in\underline J}y_i^2=0
\]
sur laquelle le groupe $V_{\underline J}$ \inserted{$\simeq\mathbb F_2^{\underline J}/\mathbb F_2$} opère \deleted{\unsure{par}} par changements de signes sur les $x_i$, \unsure{a} comme quotient $\Sigma_{\underline J}$, défini par l'équation homogène
\[
(13)\qquad\sum x_i=0.
\]
\margin{description valable loc. (ét.) ; sinon, il faut prendre dans $\mathbb P(p_*(\mathcal O_{\underline J}))$, \deleted{\unsure{de}} la conique d'équation $Q=0$, où $Q(u\deleted{\unsure{u}})=\operatorname{Tr}u^2$ (NB $p:\underline J\longrightarrow S$ le morphisme structural).}
\margin{description valable loc. ét. -- la description est par l'hyperplan (13) de $\mathbb P(p_*(\mathcal O_{\underline J}))$, d'équation $\unsure{L(u)}=0$, où $L(u)=\operatorname{Tr}u$.}

Donc La donnée de $(X,\underline I)$ équivaut à celle de \deleted{\unsure{$\Sigma$}} $\underline J$ i.e. de $\Sigma_{\underline J}$, d'un torseur $\underline K$ sous $V_{\underline J}$ i.e. d'un revêtement $X$ de $\Sigma_{\underline J}$ \deleted{ramifié} de groupe $V_{\underline J}$ (admissible au sens qu'il définit l'\uline{identité} $V_{\underline J}\longrightarrow V_{\underline J}$), \uline{et} d'une section $\unsure{t|x}$ de $\Sigma^*_{\underline J}$.

Quand $\underline I$ est \inserted{fixé et} de la forme $I_S$, avec $I$ ens. à quatre éléments fixe (à fins d'épinglage), on a $\underline J\simeq J_S$, et $\Sigma^*_{\underline J}=\bar\Sigma_J$ est aussi fixé. La donnée de $(X,I_S\hookrightarrow X)$ équivaut alors à celle de la section $x$ de $\Sigma^*_J$, servant de ``pt référence'', le revêtement $X$ de $\bar\Sigma^*_J$ de groupe $(V_{\underline J})_S=(V)_{(J_S)}=\underline V$ étant
\margin{[L'hypothèse que $\underline I$ soit constant est en fait superflue]}

\margin{Relations entre $\underline I,\underline K,x$ :
\[
(14)\quad\underline I\simeq\underline K\wedge_{\underline V}\bigl(\unsure{\widetilde\Sigma}_{\underline J,2}^*|x\bigr).
\]
Ainsi, pour $\underline J$ et $x\in\Gamma(\Sigma^*_{\underline J}/S)$ fixés, $\underline I$ et $\underline K$ (ou encore, $\underline I$ et $X$) se déterminent mutuellement\ldots On a trois revêtements étales de degré $4$ de $S$, $\underline I$, $\underline J\amalg S$, et $\underline K$\ldots}
